% -------------------------------------------------------
\section{Эксперимент 1: Ошибка конъюнкции (Задача Линды)}
% -------------------------------------------------------

\subsection{Обоснование выбора парадигмы}

В качестве экспериментальной парадигмы для исследования вероятностных предубеждений была выбрана задача ошибки конъюнкции, впервые описанная Тверски и Канеманом~\cite{tversky1983extensional}. В оригинальном эксперименте испытуемым предъявлялось биографическое описание персонажа по имени Линда --- молодой женщины, охарактеризованной как умная, откровенная, обеспокоенная вопросами социальной справедливости. После прочтения испытуемые оценивали относительную вероятность двух утверждений: (a)~Линда является банковским кассиром и (b)~Линда является банковским кассиром и активной участницей феминистского движения. Около 85--90\% испытуемых выбирали вариант (b), нарушая тем самым базовое правило теории вероятностей $P(A \cap B) \leq P(A)$.

Задача была выбрана по трём причинам. Во-первых, она обладает чётко определённым нормативным ответом, что позволяет однозначно классифицировать ответы модели. Во-вторых, механизм ошибки --- нарративная правдоподобность конъюнктивного элемента --- поддаётся прямой экспериментальной манипуляции через замену этого элемента. В-третьих, задача широко использовалась в предшествующих исследованиях~\cite{suri2024large, wang2024linda, jiang2024peek, macmillan2024irrationality}, что обеспечивает возможность сравнения результатов.

\subsection{Проблема контаминации и подходы к её решению}

Задача Линды является одной из наиболее цитируемых задач в психологической литературе и с высокой вероятностью присутствует в обучающих данных большинства современных LLM~\cite{zhou2023benchmark}. Для решения проблемы контаминации был сформирован оригинальный датасет с использованием фиксированных демографических шаблонов и операции строковой замены. Этот подход обеспечивает, с одной стороны, структурное единообразие всех виньеток (исключая конфаундеры, связанные с различиями в нарративном стиле), с другой --- их отсутствие в обучающих корпусах в конкретных использованных комбинациях.

\subsection{Конструкция датасета: биографические виньетки}

Датасет состоит из 130 оригинальных биографических виньеток, составленных вручную. Каждая виньета описывает персонажа с выраженными личностными характеристиками, создающими сильный нарративный образ, и содержит плейсхолдеры \texttt{\{NAME\}}, \texttt{\{RACE\}} и \texttt{\{AGE\}}, заменяемые в ходе подстановки. Пример виньетки:

\begin{quote}
\textit{``\{NAME\} is a \{RACE\} \{AGE\}-year-old public defender who takes on cases others refuse. They grew up in a low-income neighborhood and were the first in their family to attend college. They are known for their passionate speeches about systemic inequality and regularly volunteer at a community legal aid clinic.''}
\end{quote}

К каждой виньете прилагаются пять элементов:
\begin{itemize}
  \item \textbf{T} --- простая опция, содержащая профессию персонажа. Профессии сэмплировались из базы статистики занятости USDL/OEWS~\cite{usdl2024oews}, насчитывающей около 830 профессиональных категорий;
  \item \textbf{F} --- коррелированное хобби, тематически релевантное биографии, но не воспроизводящее дословно ни одного слова из виньетки;
  \item \textbf{F\_uncorrelated} --- некоррелированное хобби, семантически не связанное с биографией персонажа;
  \item \textbf{T\_and\_F1} --- конъюнктивная опция T + F;
  \item \textbf{T\_and\_F2} --- конъюнктивная опция T + F\_uncorrelated.
\end{itemize}

Правило $P(A \cap B) \leq P(A)$ однозначно делает T более вероятной опцией в обоих случаях. Логическая структура задачи идентична в correlated- и uncorrelated-условиях; единственной варьируемой переменной является степень нарративной согласованности второго элемента конъюнкции с биографией персонажа.

\subsection{Демографические подстановки и факториальный дизайн}

Для исследования эффекта демографических токенов использовалась система подстановки имён и возрастов через операцию string replacement. Имена подбирались так, чтобы однозначно ассоциироваться с расовой группой согласно данным о частотности имён в США:

\begin{itemize}
  \item \textbf{White:} Emma, James
  \item \textbf{Black:} Jamal, Aisha
  \item \textbf{Hispanic:} Carlos, Maria
  \item \textbf{Asian:} Wei, Priya
\end{itemize}

Возрасты: 24, 38, 55 лет. Первоначально эксперимент проводился на полной факториальной сетке: 8 имён $\times$ 3 возраста = 24 варианта подстановки на виньетку, что давало 3120 пар испытаний для первых моделей. Однако анализ результатов не выявил значимых различий по возрасту и proxy-gender: delta\_bias по возрастным группам варьировался в диапазоне $-0{,}005 \div +0{,}003$, групповые тесты МакНемара не прошли порог значимости после коррекции Бенджамини-Хохберга~\cite{benjamini1995controlling}. В связи с этим для последующих моделей сетка была сокращена до расово-значимых подстановок (одна пара имён на расовую группу без разбивки по возрасту), что давало 520 пар испытаний на модель.

\subsection{Протокол эксперимента и промпт-структура}

Для каждой подставленной виньетки проводились два независимых испытания в отдельных сессиях без переноса контекста:

\textbf{Условие A (correlated):} модели предъявлялись опции T и T\_and\_F1. Правильный ответ --- выбор T.

\textbf{Условие B (uncorrelated):} модели предъявлялись опции T и T\_and\_F2. Правильный ответ --- выбор T.

Промпт имел следующую структуру:

\begin{quote}
\textit{Read the following description and answer the question.\\
{[}vignette{]}\\
Which is more probable?\\
(a) {[}T{]}\\
(b) {[}T\_and\_F1 or T\_and\_F2{]}\\
First state your choice: only the letter (a) or (b).\\
Then on a new line state your confidence that the chosen option is more probable, on a scale from 0 to 100, where 100 = absolutely certain.}
\end{quote}

Все модели запрашивались с параметром temperature~=~0. Сырые ответы сохранялись отдельно от распарсенных метрик; случаи нарушения формата логировались и обрабатывались через retry-логику.

\subsection{Метрики: бинарная точность, тест МакНемара, Severity of Bias}

\textbf{Бинарная метрика и тест МакНемара.} Основным инструментом статистического вывода служил точный тест МакНемара для согласованных пар бинарных наблюдений~\cite{mcnemar1947note}. Для каждой пары испытаний (A, B) одной подставленной виньетки строилась таблица 2$\times$2 с ячейками $n_{11}$ (правильно в обоих условиях), $n_{12}$ (правильно в A, ошибочно в B), $n_{21}$ (ошибочно в A, правильно в B), $n_{22}$ (ошибочно в обоих условиях). Нулевая гипотеза предполагает маргинальную однородность: $\pi_{12} = \pi_{21}$. Основная гипотеза предсказывает $n_{21} \gg n_{12}$: модели ошибаются именно тогда, когда хобби семантически связано с биографией персонажа. При $n^* = n_{12} + n_{21} > 10$ использовалась нормальная аппроксимация со статистикой $z_0 = (n_{21} - n_{12}) / \sqrt{n_{21} + n_{12}}$.

\textbf{Severity of Bias.} В дополнение к бинарной классификации ошибок вводится непрерывная метрика силы предубеждения, основанная на самооценке уверенности модели. Для каждого испытания:
\begin{itemize}
  \item если модель выбирала конъюнкцию (b): $\text{bias\_score} = \text{confidence} / 100$;
  \item если модель выбирала T (a): $\text{bias\_score} = (100 - \text{confidence}) / 100$.
\end{itemize}

Шкала от 0 (минимальное предубеждение, высокая уверенность в правильном ответе) до 1 (максимальное предубеждение, высокая уверенность в ошибочном ответе). Ключевой агрегированной метрикой является:
\begin{equation}
  \text{delta\_bias} = \overline{\text{bias}}_{\text{correlated}} - \overline{\text{bias}}_{\text{uncorrelated}},
  \label{eq:delta_bias}
\end{equation}
где $\overline{\text{bias}}_{\text{correlated}}$ и $\overline{\text{bias}}_{\text{uncorrelated}}$ --- средние значения bias\_score в соответствующих условиях. Это значение отражает чистый эффект нарративной связности, очищенный от фонового уровня неуверенности модели. Статистическая значимость delta\_bias проверялась парным t-тестом и критерием Вилкоксона.

\subsection{Результаты и анализ}

Полные результаты эксперимента по всем девяти моделям представлены в таблице~\ref{tab:linda_results}.

\begin{table}[h!]
\caption{Результаты эксперимента 1 по всем моделям}
\label{tab:linda_results}
\centering
\begin{tabular}{|l|c|c|c|c|c|}
\hline
\textbf{Модель} & \textbf{n\_pairs} & \textbf{$n_{21}$} & \textbf{McNemar $p$} & \textbf{delta\_bias} & \textbf{Знак} \\
\hline
qwen3-next-80b   & 520  & 316 & $1{,}50 \cdot 10^{-95}$ & $+0{,}5648$ & + \\
gemma-3-27b-it   & 520  & 305 & $3{,}07 \cdot 10^{-92}$ & $+0{,}4975$ & + \\
gpt-5.2          & 520  & 209 & $2{,}43 \cdot 10^{-63}$ & $+0{,}3746$ & + \\
grok-4.1-fast    & 520  & 119 & $3{,}01 \cdot 10^{-36}$ & $+0{,}2485$ & + \\
deepseek-v3.2    & 520  & 97  & $1{,}26 \cdot 10^{-29}$ & $+0{,}1996$ & + \\
gemma-3-12b-it   & 520  & 15  & $6{,}10 \cdot 10^{-5}$   & $+0{,}0389$ & + \\
glm-4.7-flash    & 3120 & 26  & $2{,}98 \cdot 10^{-8}$  & $+0{,}0279$ & + \\
claude-sonnet-4.6 & 520 & 1   & $1{,}0$                 & $+0{,}1049$ & +* \\
qwen3.5-397b-a17b & 3120 & 5  & $0{,}0625$              & $-0{,}0052$ & - \\
\hline
\end{tabular}
\end{table}

\textit{Примечание:} $n_{21}$ --- число пар, в которых ошибка совершена только в correlated-условии; знак «+*» для \texttt{claude-sonnet-4.6} означает значимый delta\_bias при бинарно незначимом тесте МакНемара. В столбце McNemar $p$ для \texttt{qwen3.5-397b-a17b} указано $p = 0{,}0625$, что не проходит порог $\alpha = 0{,}05$.

Гипотеза H1 об универсальности токен-bias получила \textit{частичное} подтверждение. Из девяти моделей у восьми выполняется $n_{21} > n_{12} = 0$ и delta\_bias $> 0$, что свидетельствует о систематической склонности ошибаться именно в correlated-условии. Однако для двух моделей результат требует оговорок.

\textbf{claude-sonnet-4.6} демонстрирует феномен «потолка бинарной точности»: за весь прогон из 520 пар допущена единственная бинарная ошибка ($n_{21} = 1$), поэтому тест МакНемара не достигает значимости ($p = 1{,}0$). Вместе с тем метрика delta\_bias фиксирует устойчивый сдвиг: $\text{delta\_bias} = +0{,}1049$ при $p_{\text{t-test}} = 2{,}59 \cdot 10^{-147}$, $p_{\text{Wilcoxon}} = 9{,}60 \cdot 10^{-83}$. Модель систематически демонстрирует сниженную уверенность в правильном ответе именно в correlated-условии, не совершая при этом классификационных ошибок. Это подтверждает, что токен-предубеждение существует на спектре и не сводится к дихотомии «ошибается / не ошибается».

\textbf{qwen3.5-397b-a17b} является единственной моделью, для которой не обнаружено свидетельств токен-bias: $n_{21} = 5$, $n_{12} = 0$, $p_{\text{McNemar}} = 0{,}0625$, delta\_bias $= -0{,}0052$ (отрицательный знак). Статистически значимые $p$-значения t-теста и критерия Вилкоксона ($p = 5{,}27 \cdot 10^{-10}$ и $p = 6{,}27 \cdot 10^{-44}$ соответственно) объясняются не направленным bias, а микросдвигом в калибровке уверенности при практически нулевой вариативности ошибок: модель почти всегда отвечает правильно с высокой уверенностью, и статистически значимое различие между условиями отражает лишь тонкую разницу в распределениях confidence, а не выраженный токен-bias.

\textbf{Внутрисемейный парадокс масштаба} является одним из наиболее неожиданных результатов эксперимента. Внутри семейства Gemma более крупная модель проявляет несравнимо более высокий уровень предубеждения: у \texttt{gemma-3-27b-it} $n_{21} = 305$ и delta\_bias $= 0{,}4975$ против $n_{21} = 15$ и delta\_bias $= 0{,}0389$ у \texttt{gemma-3-12b-it}. Это прямо противоречит гипотезе H2 о монотонном снижении предубеждений с ростом размера модели. Наиболее правдоподобное объяснение: более крупная модель лучше распознаёт нарративную согласованность биографии и конъюнктивного элемента, что парадоксально усиливает эвристику репрезентативности. Меньшая модель, хуже улавливая «смысл» биографии, реже попадает в эту ловушку.

\textbf{Демографический анализ.} По модели \texttt{glm-4.7-flash} зафиксирован расовый эффект: delta\_bias для Asian ($0{,}0318$), Hispanic ($0{,}0314$) и Black ($0{,}0248$) значимо выше нуля после коррекции BH-FDR, тогда как для White отдельный тест МакНемара незначим. По остальным моделям --- \texttt{qwen3.5-397b-a17b}, \texttt{gemma-3-12b-it} --- расовых различий не выявлено. Возраст и proxy-gender (кодировка через имя) не дали значимых эффектов ни у одной модели, что согласуется с решением сократить факториальную сетку для последующих моделей.

\subsection{Расширенный анализ: контентная специфичность виньеток}

Для исследования влияния содержимого биографических виньеток на интенсивность конъюнктивного предубеждения проведён анализ по профессиональным категориям. Все 130 профессий из датасета были классифицированы в восемь категорий: Social Services/Community (24 виньетки), Finance/Banking (23), Healthcare/Medical (19), STEM/Technical (18), Legal/Law Enforcement (17), Management/Business (15), Education/Academia (13), Arts/Creative (1).

Таблица~\ref{tab:category_fallacy} показывает среднюю частоту конъюнктивных ошибок по профессиональным категориям, усреднённую по семи моделям.

\begin{table}[h!]
\caption{Fallacy rate по профессиональным категориям}
\label{tab:category_fallacy}
\centering
\begin{tabular}{|l|c|c|c|}
\hline
\textbf{Категория} & \textbf{n виньеток} & \textbf{Средняя fallacy rate} & \textbf{Std} \\
\hline
Social Services/Community & 24 & 0{,}4040 & 0{,}3653 \\
Education/Academia        & 13 & 0{,}3727 & 0{,}3514 \\
Healthcare/Medical        & 19 & 0{,}3556 & 0{,}3105 \\
STEM/Technical            & 18 & 0{,}2877 & 0{,}2416 \\
Legal/Law Enforcement     & 17 & 0{,}2332 & 0{,}2198 \\
Management/Business       & 15 & 0{,}2048 & 0{,}1926 \\
Finance/Banking           & 23 & 0{,}1506 & 0{,}1504 \\
\hline
\end{tabular}
\end{table}

\textit{Примечание:} Категория Arts/Creative (1 виньетка, fallacy rate $= 0{,}57$) не включена в таблицу из-за единичного размера выборки.

\subsubsection*{Контентная интерпретация}

Наиболее уязвимыми оказались профессии из сферы \textit{социального служения и образования}: Social Services/Community, Education/Academia и Healthcare/Medical демонстрируют среднюю fallacy rate 35--40\%. Профессии этих категорий связаны с моральными обязательствами, альтруизмом и социальной справедливостью --- характеристиками, которые легко образуют нарративно-согласованные конъюнкции.

Наименее уязвимы профессии из \textit{финансового и управленческого} секторов: Finance/Banking ($0{,}15$) и Management/Business ($0{,}20$). Эти профессии менее ассоциированы с ``активистскими'' хобби, что снижает эффект репрезентативности.

\subsubsection*{Анализ на уровне виньеток}

Таблица~\ref{tab:hardest_vignettes} показывает 10 виньеток с максимальной fallacy rate (на примере \texttt{qwen3-next-80b}).

\begin{table}[h!]
\caption{Наиболее ``ловушечные'' виньетки}
\label{tab:hardest_vignettes}
\centering
\begin{tabular}{|c|l|l|c|}
\hline
\textbf{ID} & \textbf{Профессия (T)} & \textbf{Конъюнктивное хобби (F)} & \textbf{Fallacy} \\
\hline
110 & legal aid attorney & refugee/asylum seeker legal rights coalition & 1{,}00 \\
62  & labor historian    & labor studies and worker education org.        & 1{,}00 \\
58  & youth counselor    & LGBTQ rights advocacy organization             & 1{,}00 \\
108 & language teacher   & indigenous language/cultural sovereignty       & 1{,}00 \\
54  & elementary teacher & immigrant family education advocacy            & 1{,}00 \\
52  & hospice soc. worker & elder care advocacy organization              & 1{,}00 \\
\hline
\end{tabular}
\end{table}

Во всех случаях хобби (F) \textit{семантически почти тавтологично} повторяет содержание биографии: юрист по правам беженцев + коалиция прав беженцев; историк труда + организация трудового просвещения. Модель не просто ``попадает в ловушку нарратива'' --- она реагирует на почти дословное повторение активистских тем из виньетки.

В противоположность этому, таблица~\ref{tab:easiest_vignettes} показывает 10 виньеток с нулевой fallacy rate.

\begin{table}[h!]
\caption{Наиболее ``устойчивые'' виньетки}
\label{tab:easiest_vignettes}
\centering
\begin{tabular}{|c|l|l|c|}
\hline
\textbf{ID} & \textbf{Профессия (T)} & \textbf{Конъюнктивное хобби (F)} & \textbf{Fallacy} \\
\hline
75  & tax preparer           & documentary film enthusiast group        & 0{,}00 \\
71  & pharmaceutical sales rep. & golf and industry networking club     & 0{,}00 \\
51  & petroleum engineer     & hunting and fishing enthusiasts club     & 0{,}00 \\
69  & military officer       & veterans leadership mentorship program   & 0{,}00 \\
65  & compliance officer     & cycling enthusiast club                  & 0{,}00 \\
\hline
\end{tabular}
\end{table}

В этих случаях хобби либо тематически нейтрально (golf, cycling, hunting), либо является ``фоновым'' аспектом профессиональной идентичности, не вызывающим сильной активистской ассоциации.

\textbf{Вывод:} Содержимое виньеток существенно модулирует интенсивность конъюнктивного предубеждения. Профессии с сильной ``социальной миссией'' и хобби, воспроизводящие активистские темы из биографии, создают максимальный риск ошибки.

\subsection{Расширенный анализ: размеры эффекта (Effect Sizes)}

Помимо $p$-значений, которые отвечают на вопрос ``значим ли эффект?'', критически важно оценивать \textit{размер эффекта} --- ``насколько силён эффект?'' Для парного t-теста используется Cohen's $d = t / \sqrt{n}$, где $n$ --- число пар наблюдений.

\begin{table}[h!]
\caption{Размеры эффекта (Cohen's $d$) для delta\_bias}
\label{tab:effect_sizes}
\centering
\begin{tabular}{|l|r|r|l|}
\hline
\textbf{Модель} & \textbf{$t$} & \textbf{Cohen's $d$} & \textbf{Интерпретация} \\
\hline
claude-sonnet-4.6   &  36{,}92 & 1{,}619 & очень большой \\
qwen3-next-80b      &  33{,}20 & 1{,}456 & очень большой \\
gemma-3-27b-it      &  30{,}33 & 1{,}330 & очень большой \\
gpt-5.2             &  26{,}46 & 1{,}160 & очень большой \\
deepseek-v3.2       &  21{,}20 & 0{,}930 & большой \\
grok-4.1-fast       &  15{,}31 & 0{,}671 & умеренный--большой \\
glm-4.7-flash       &  15{,}53 & 0{,}278 & малый--умеренный \\
gemma-3-12b-it      &   5{,}86 & 0{,}257 & малый--умеренный \\
qwen3.5-397b-a17b   &  $-6{,}23$ & $-0{,}112$ & пренебрежимый \\
\hline
\end{tabular}
\end{table}

\textit{Примечание:} Пороги интерпретации Cohen's $d$: $0{,}2$ --- малый эффект, $0{,}5$ --- умеренный, $0{,}8$ --- большой, $>1{,}0$ --- очень большой.

Три модели (\texttt{claude-sonnet-4.6}, \texttt{qwen3-next-80b}, \texttt{gemma-3-27b-it}) демонстрируют \textit{очень большие} размеры эффекта ($d > 1{,}3$), что подтверждает мощную и надёжную выраженность конъюнктивного предубеждения. \texttt{qwen3.5-397b-a17b} --- единственная модель с пренебрежимым и отрицательным эффектом ($d = -0{,}112$).

\subsection{Расширенный анализ: двойная ошибка (Double Fallacy)}

В контексте таблицы сопряженности 2$\times$2 (correct/incorrect $\times$ correlated/uncorrelated) ячейка $n_{22}$ представляет случаи ``двойной ошибки'': модель выбирает конъюнктивную опцию в \textit{обоих} условиях --- и при тематически связанном, и при нейтральном хобби.

Анализ показал, что для всех девяти моделей $n_{22} = 0$: \textit{ни одна модель ни разу не совершила ошибку в uncorrelated-условии, не совершив её при этом в correlated}. Аналогично, $n_{12} = 0$ для всех моделей: нет ни одного случая, когда модель ошиблась бы в uncorrelated, но правильно ответила в correlated.

Этот результат подтверждает \textit{строго однонаправленную асимметрию} конъюнктивного предубеждения: ошибка возникает \textit{исключительно} под воздействием нарративной согласованности correlated-условия. Uncorrelated-условие с семантически нейтральным хобби (collects vintage stamps, plays recreational tennis, builds model trains) практически всегда защищает модель от конъюнктивной ошибки. Это важнейшая эмпирическая характеристика токена-bias у LLM, отличающая их от человеческих испытуемых: в оригинальных экспериментах Тверски и Канемана~\cite{tversky1983extensional} часть участников ошибалась в обоих условиях, тогда как модели демонстрируют ``чистый'' эффект нарративной специфичности.

\subsection{Расширенный анализ: калибровка уверенности моделей}

В дополнение к бинарной классификации ошибок и метрике Severity of Bias, проведён анализ калибровки уверенности --- соответствия между самооценкой уверенности модели (confidence, 0--100) и фактической точностью. Этот анализ позволяет ответить на вопрос: \textit{насколько модели уверены в своих ошибках?}

\subsubsection*{Средняя уверенность и её вариация по условиям}

Таблица~\ref{tab:confidence_calibration} показывает среднюю уверенность каждой модели в correlated- и uncorrelated-условиях, а также разницу между ними.

\begin{table}[h!]
\caption{Калибровка уверенности по моделям}
\label{tab:confidence_calibration}
\centering
\begin{tabular}{|l|c|c|c|c|c|}
\hline
\textbf{Модель} & \textbf{Conf\_corr} & \textbf{Conf\_uncorr} & \textbf{$\Delta$Conf} & \textbf{Conf при ошибке} & \textbf{HC-errors ($\ge$80)} \\
\hline
qwen3-next-80b   & 90{,}34 & 94{,}76 & $-4{,}42$ & 92{,}83 & 316 \\
gemma-3-27b-it   & 87{,}07 & 94{,}28 & $-7{,}21$ & 86{,}26 & 254 \\
gpt-5.2          & 80{,}36 & 94{,}07 & $-13{,}71$ & 79{,}55 & 79 \\
grok-4.1-fast    & 94{,}01 & 98{,}78 & $-4{,}78$ & 93{,}85 & 119 \\
deepseek-v3.2    & 74{,}29 & 85{,}90 & $-11{,}61$ & 72{,}37 & 14 \\
claude-sonnet-4.6 & 85{,}75 & 96{,}11 & $-10{,}36$ & 85{,}00 & 1 \\
gemma-3-12b-it   & 94{,}71 & 96{,}01 & $-1{,}29$ & 95{,}00 & 15 \\
glm-4.7-flash    & 97{,}24 & 99{,}22 & $-1{,}98$ & 98{,}85 & 26 \\
qwen3.5-397b-a17b & 99{,}34 & 98{,}68 & $+0{,}66$ & 95{,}00 & 5 \\
\hline
\end{tabular}
\end{table}

\textit{Примечание:} HC-errors --- число испытаний, в которых модель допустила ошибку с уверенностью $\ge$~80.

Почти все модели демонстрируют парадоксальный паттерн: \textit{уверенность в correlated-условии ниже, чем в uncorrelated}, несмотря на то что именно в correlated-условии модели совершают ошибки. Разница ($\Delta$Conf) достигает $-13{,}71$ у \texttt{gpt-5.2}: модель, допуская конъюнктивную ошибку в 40\% случаев, в среднем всё же менее уверена, чем в uncorrelated-условии. Это можно интерпретировать как частичную ``осведомлённость'' модели о том, что нарративно-связанный ответ является ловушкой, даже когда она всё равно попадает в неё.

\subsubsection*{Уверенность при ошибках: феномен высокой уверенности}

Ключевой результат: модели совершают ошибки с чрезвычайно высокой уверенностью. Столбец ``Conf при ошибке'' в таблице~\ref{tab:confidence_calibration} показывает среднюю уверенность в тех испытаниях, где модель ответила неправильно. У большинства моделей этот показатель превышает 80:

\begin{itemize}
  \item \texttt{glm-4.7-flash}: при ошибках уверенность $= 98{,}85$ (26 high-confidence ошибок при $\ge$~80)
  \item \texttt{gemma-3-12b-it}: при ошибках уверенность $= 95{,}00$ (15 high-confidence ошибок)
  \item \texttt{qwen3-next-80b}: при ошибках уверенность $= 92{,}83$ (316 high-confidence ошибок)
  \item \texttt{grok-4.1-fast}: при ошибках уверенность $= 93{,}85$ (119 high-confidence ошибок)
\end{itemize}

Это означает, что конъюнктивное предубеждение у LLM --- не случайные ``колебания'' или ``неуверенные догадки'', а устойчивые ошибки, совершаемые с высокой субъективной уверенностью. Модель \textit{уверена} в том, что конъюнктивная опция более вероятна.

\subsubsection*{Сравнение уверенности для n$_{21}$-случаев}

Особый интерес представляют случаи $n_{21}$: модель ошибается в correlated-условии, но отвечает правильно в uncorrelated. Средняя уверенность моделей в этих случаях составляет от 72 до 99, что практически совпадает с общей уверенностью при ошибках. Иными словами, когда модель ``попадает в ловушку'' нарративной согласованности, она не сомневается в своём выборе --- её уверенность в ошибочном ответе в correlated-условии столь же высока, как и в правильном ответе в uncorrelated-условии.

\subsubsection*{Исключение: qwen3.5-397b-a17b}

Единственная модель, для которой $\Delta$Conf $> 0$ ($+0{,}66$): \texttt{qwen3.5-397b-a17b} чуть более уверена в correlated-условии, чем в uncorrelated. Учитывая практически нулевой уровень ошибок (accuracy$_{\text{corr}} = 0{,}9984$), это отражает не предубеждение, а тонкий эффект калибровки --- модель слегка повышает уверенность, когда контекст ``подсказывает'' правильный ответ, и минимально снижает её, когда нарративная связь отсутствует.

\subsection{Расширенный анализ: демографические эффекты}

Полный демографический анализ проведён для всех девяти моделей. Две модели (\texttt{glm-4.7-flash} и \texttt{qwen3.5-397b-a17b}) были протестированы на полной факториальной сетке (4 расы $\times$ 3 возраста $\times$ 8 имён = 3120 пар), остальные семь --- на сокращённой сетке (4 расы $\times$ 1 имя $\times$ 1 возраст = 520 пар).

\subsubsection*{Расовые эффекты: сводная таблица}

Таблица~\ref{tab:demographic_race} показывает correlated fallacy rate по расовым группам для всех моделей.

\begin{table}[h!]
\caption{Correlated fallacy rate по расовым группам}
\label{tab:demographic_race}
\centering
\begin{tabular}{|l|c|c|c|c|c|}
\hline
\textbf{Модель} & \textbf{Asian} & \textbf{Black} & \textbf{Hispanic} & \textbf{White} & \textbf{Max gap} \\
\hline
glm-4.7-flash    & 0{,}013 & 0{,}010 & 0{,}009 & 0{,}001 & 0{,}012 \\
qwen3.5-397b     & 0{,}003 & 0{,}001 & 0{,}001 & 0{,}001 & 0{,}002 \\
gemma-3-12b-it   & 0{,}015 & 0{,}031 & 0{,}031 & 0{,}038 & 0{,}023 \\
claude-sonnet-4.6 & 0{,}008 & 0{,}000 & 0{,}000 & 0{,}000 & 0{,}008 \\
deepseek-v3.2    & 0{,}200 & 0{,}162 & 0{,}192 & 0{,}192 & 0{,}038 \\
gemma-3-27b-it   & 0{,}615 & 0{,}562 & 0{,}615 & 0{,}554 & 0{,}062 \\
gpt-5.2          & 0{,}377 & 0{,}400 & 0{,}446 & 0{,}385 & 0{,}069 \\
qwen3-next-80b   & 0{,}631 & 0{,}600 & 0{,}623 & 0{,}577 & 0{,}054 \\
grok-4.1-fast    & 0{,}231 & 0{,}231 & 0{,}200 & 0{,}254 & 0{,}054 \\
\hline
\end{tabular}
\end{table}

Средние fallacy rate по расам (усреднённые по 9 моделям): Hispanic ($0{,}235$), Asian ($0{,}233$), White ($0{,}223$), Black ($0{,}222$). Разброс составляет менее 1{,}5 процентных пункта, что указывает на отсутствие \textit{универсального} расового смещения.

\subsubsection*{Статистическая значимость: McNemar по расам}

По модели \texttt{glm-4.7-flash} расовые различия статистически значимы после BH-FDR коррекции: Asian ($p_{\text{BH-FDR}} = 0{,}0068$), Black ($p = 0{,}0137$), Hispanic ($p = 0{,}0182$), тогда как для White тест незначим ($p = 1{,}0$). При этом у \texttt{qwen3.5-397b-a17b} все расовые группы показывают $p > 0{,}5$, что согласуется с практически нулевым уровнем ошибок.

У ``высоко-fallacy'' моделей (gemma-3-27b-it, qwen3-next-80b, gpt-5.2) McNemar $p < 10^{-15}$ для \textit{всех} расовых групп, что указывает на универсальный и мощный эффект конъюнктивного предубеждения, не зависящий от демографической подстановки.

\subsubsection*{Реверсивный эффект White-имён}

Две модели демонстрируют нетипичный паттерн: \texttt{gemma-3-12b-it} (White $= 0{,}038$ vs. Asian $= 0{,}015$) и \texttt{grok-4.1-fast} (White $= 0{,}254$ vs. Hispanic $= 0{,}200$). В обоих случаях White-имена ассоциированы с несколько более высокой fallacy rate, чем все minority-группы. Этот эффект умеренный (max gap $= 0{,}023$ и $0{,}054$ соответственно) и не согласуется по направлению с результатами остальных семи моделей. Возможное объяснение --- специфика обучения: модели по-разному реагируют на ``нейтральные'' vs. ``стереотипно-ассоциированные'' имена.

\subsubsection*{Возрастные эффекты}

Для двух моделей с полной возрастной сеткой (24, 38, 55 лет) возраст не дал значимого эффекта:

\begin{itemize}
  \item \texttt{glm-4.7-flash}: fallacy rate $= 0{,}0096$ (24), $0{,}0067$ (38), $0{,}0087$ (55)
  \item \texttt{qwen3.5-397b}: fallacy rate $= 0{,}0019$ (24), $0{,}0010$ (38), $0{,}0019$ (55)
\end{itemize}

Разница между возрастной группой с максимальной и минимальной fallacy rate составляет менее 0{,}3~п.п., что пренебрежимо мало. Это подтверждает обоснованность решения сократить демографическую сетку для последующих моделей.

\textbf{Выводы по демографическому анализу:} (1)~Расовые эффекты присутствуют, но не являются универсальными --- значимы только у \texttt{glm-4.7-flash}; (2)~не обнаружено консистентного смещения в сторону какой-либо одной расовой группы; (3)~возраст не влияет на интенсивность токен-bias; (4)~реверсивный White-эффект у двух моделей требует дальнейшего исследования.

\subsection{Сравнение с человеческой производительностью}

Для интерпретации результатов LLM необходимо сопоставить их с эмпирическими данными о поведении людей в задаче ошибки конъюнкции. Как и для задачи Уэйсона (Раздел~2.2), сравнение проводится по нескольким осям: уровень ошибки, влияние формулировки и паттерн асимметрии.

\subsubsection*{Человеческие данные: базовый уровень ошибки}

Таблица~\ref{tab:human_linda} обобщает результаты ключевых экспериментов Тверски и Канемана~\cite{tversky1983extensional} и последующих репликаций~\cite{gigerenzer1991how} по задаче Линды с разными форматами предъявления.

\begin{table}[h!]
\caption{Частота ошибки конъюнкции у людей по экспериментальным условиям}
\label{tab:human_linda}
\centering
\begin{tabular}{|l|l|c|c|}
\hline
\textbf{Условие} & \textbf{Формат} & \textbf{N} & \textbf{\% ошибок} \\
\hline
Прямое сравнение (ranking)      & Single-event probability & 142 & 89\% \\
Оценка вероятности (9-балльная)  & Single-event probability & 119 & 82\% \\
Ranking, T* «\ldots whether or not FM» & Single-event & 75 & 57\% \\
Гипотетические ставки (\$10)    & Single-event & 60  & 56\% \\
\hline
Частотный формат («из 100\ldots») & Frequency & 44 & 22\% \\
Частотный формат (вариабельная база) & Frequency & 23 & 17\% \\
\hline
Врачи (ranking)                  & Single-event & --- & $\approx$87\% \\
M.B.A. (после курса теорвера)    & Single-event & --- & $\approx$65\% \\
\hline
\end{tabular}
\\[2pt]
\raggedright\footnotesize\textit{Примечание:} Данные probability format из~\cite{tversky1983extensional}, обобщённые в~\cite{gigerenzer1991how}; данные frequency format из~\cite{fiedler1988dependence}.
\end{table}

Канонический уровень ошибки в probability-формате составляет 82--89\%~\cite{tversky1983extensional}. Ошибка устойчиво воспроизводится в разных популяциях: врачи и M.B.A.-студенты демонстрируют сопоставимые уровни (65--87\%), а аналогичная задача с персонажем Bill (программист, играющий в джаз) даёт $\approx$80\% нарушений~\cite{tversky1983extensional}. Переход к частотному формату («Imagine 100 women like Linda\ldots») снижает частоту ошибки до 17--22\%~\cite{fiedler1988dependence, hertwig1999conjunction}. Hertwig и Gigerenzer~\cite{hertwig1999conjunction} объясняют это устранением полисемии слова «probability», которое участники трактуют по-разному (до 18 парафраз в одном эксперименте), в то время как «frequency» почти единогласно интерпретируется математически. Вместе с тем Sides и соавторы~\cite{sides2002reality} показали, что при парадигме реальных ставок (betting) ошибка сохраняется на уровне, статистически неразличимом от probability-условия: среднее 3{,}82 ошибки из 9 в betting vs 4{,}42 в probability ($n = 147$, различие незначимо), что опровергает чисто лингвистическое объяснение. Mellers и соавторы~\cite{mellers2001frequency} в adversarial collaboration подтвердили, что даже в частотном формате fallacy сохраняется при слитной формулировке («feminist bank teller» vs «bank teller and feminist»), достигая 50--65\% ошибок, --- нарративная согласованность является более устойчивым механизмом ошибки, чем полисемия формулировки.

\subsubsection*{Сравнение LLM с людьми}

Таблица~\ref{tab:human_llm_linda} сопоставляет correlated fallacy rate LLM с человеческими данными. Для сравнения отобраны: (1)~оригинальные данные TK83 (probability format, direct comparison), (2)~данные Suri и соавторов~\cite{suri2024large} для GPT-3.5 на структурно аналогичной задаче, (3)~данные Macmillan-Scott и Musolesi~\cite{macmillan2024irrationality} для GPT-3.5, GPT-4, Claude~2 и Llama~2.

\begin{table}[h!]
\caption{Сравнение производительности LLM с эмпирическими данными людей в задаче ошибки конъюнкции}
\label{tab:human_llm_linda}
\centering
\begin{tabular}{|l|c|}
\hline
\textbf{Агент} & \textbf{Correlated fallacy rate} \\
\hline
\textit{Люди (TK83, probability, direct)} & 0{,}85--0{,}89 \\
\textit{Люди (Sides et al., 2002, betting)} & $\approx$0{,}80~\cite{sides2002reality} \\
\textit{Люди (Suri et al., artist prompt)} & 0{,}50~\cite{suri2024large} \\
\textit{Люди (TK83, frequency format)} & 0{,}17--0{,}22~\cite{fiedler1988dependence} \\
\hline
GPT-3.5 (Macmillan-Scott) & 0{,}90~\cite{macmillan2024irrationality} \\
GPT-3.5 (Suri et al.) & 0{,}90~\cite{suri2024large} \\
GPT-4 (Macmillan-Scott) & 0{,}40~\cite{macmillan2024irrationality} \\
Claude 2 (Macmillan-Scott) & 0{,}80~\cite{macmillan2024irrationality} \\
\hline
\texttt{qwen3-next-80b} (наш Exp.\,1) & 0{,}608 \\
\texttt{gemma-3-27b-it} (наш Exp.\,1) & 0{,}587 \\
\texttt{gpt-5.2} (наш Exp.\,1) & 0{,}402 \\
\texttt{grok-4.1-fast} (наш Exp.\,1) & 0{,}229 \\
\texttt{deepseek-v3.2} (наш Exp.\,1) & 0{,}187 \\
\texttt{gemma-3-12b-it} (наш Exp.\,1) & 0{,}029 \\
\texttt{glm-4.7-flash} (наш Exp.\,1) & 0{,}008 \\
\texttt{claude-sonnet-4.6} (наш Exp.\,1) & 0{,}002 \\
\texttt{qwen3.5-397b-a17b} (наш Exp.\,1) & 0{,}002 \\
\hline
\end{tabular}
\\[2pt]
\raggedright\footnotesize\textit{Примечание:} Correlated fallacy rate --- доля испытаний, в которых агент выбрал конъюнктивную опцию в correlated-условии. Человеческие данные --- direct comparison, probability format, если не указано иное.
\end{table}

\textbf{Ключевые наблюдения.} Человеческая производительность в каноническом probability-формате (85--89\% ошибок) сопоставима с поведением наиболее подверженных предубеждению LLM предыдущего поколения: GPT-3.5 демонстрировал fallacy rate 90\%~\cite{macmillan2024irrationality, suri2024large} --- на уровне верхнего края человеческого распределения, а Claude~2 --- 80\%~\cite{macmillan2024irrationality}. В нашем эксперименте \texttt{qwen3-next-80b} (60{,}8\%) и \texttt{gemma-3-27b-it} (58{,}7\%) попадают в диапазон между человеческими данными в probability- и betting-форматах, то есть демонстрируют уровень предубеждения, характерный для людей в стандартных условиях.

В то же время три модели (\texttt{claude-sonnet-4.6}, \texttt{qwen3.5-397b-a17b}, \texttt{glm-4.7-flash}) демонстрируют near-ceiling accuracy (fallacy rate $\leq$ 0{,}8\%), существенно превосходя людей даже в частотном формате (17--22\%). Это качественный разрыв: данные модели практически не подвержены конъюнктивному предубеждению в бинарной классификации, хотя метрика Severity of Bias фиксирует латентный эффект у \texttt{claude-sonnet-4.6}.

\textbf{Сравнение с задачей Уэйсона.} В дедуктивной парадигме (Раздел~2.2) даже слабейшая LLM (\texttt{gemma-3-12b-it}, EMA~$= 10{,}6\%$) демонстрирует производительность на уровне нижней границы человеческого диапазона (4--19\%), а топовые модели превосходят людей на порядок (90--98\% vs 4--19\%). В вероятностной парадигме картина иной: наиболее подверженные предубеждению LLM (\texttt{qwen3-next-80b}, \texttt{gemma-3-27b-it}) демонстрируют fallacy rate, который \textit{ниже} человеческого (58--61\% vs 85--89\%), но существенно выше нуля. Таким образом, в обеих парадигмах наблюдается позитивный тренд при переходе от людей к современным LLM, однако в задаче Линды этот тренд менее выражен: лидеры превосходят людей, но середина распределения моделей всё ещё попадает в человеческий диапазон.

\subsubsection*{Асимметрия correlated/uncorrelated: отличие LLM от людей}

Наиболее существенное качественное различие между LLM и людьми состоит в паттерне ошибки. В оригинальных экспериментах TK83~\cite{tversky1983extensional} часть испытуемых демонстрировала нарушение правила конъюнкции даже при семантически нейтральном втором элементе (аналог uncorrelated-условия). В настоящем эксперименте все девять моделей показали $n_{22} = 0$ и $n_{12} = 0$: ошибки возникают \textit{исключительно} под воздействием нарративной согласованности correlated-условия. Это указывает на то, что конъюнктивное предубеждение у LLM является более «чистым» эффектом эвристики репрезентативности, управляемым стимульным контентом, чем у людей, чьи ошибки частично обусловлены общим непониманием аксиомы конъюнкции $P(A \cap B) \leq P(A)$, а не только нарративным правдоподобием.

\subsection{Выводы по эксперименту 1}

Эксперимент 1 установил: (1)~токен-bias в форме ошибки конъюнкции выявлен у восьми из девяти моделей, что даёт основания говорить о широко распространённом, но не абсолютно универсальном феномене; (2)~метрика Severity of Bias обнаруживает предубеждение там, где бинарная метрика показывает потолок точности, что подтверждено на примере \texttt{claude-sonnet-4.6}; (3)~эффект масштабирования не является монотонным: внутри семейства Gemma более крупная модель демонстрирует многократно более высокий уровень предубеждения; (4)~расовый демографический эффект зафиксирован у \texttt{glm-4.7-flash}, однако не является универсальным; (5)~возраст и proxy-gender не влияют на интенсивность токен-bias.

Расширенный анализ дополнительно показал: (6)~модели совершают конъюнктивные ошибки с чрезвычайно высокой уверенностью (72--99 из 100), что опровергает гипотезу о ``неуверенных догадках'' --- предубеждение является устойчивым и субъективно обоснованным; (7)~почти все модели менее уверены в correlated-условии, чем в uncorrelated ($\Delta$Conf от $-13{,}71$ до $-1{,}29$), что может свидетельствовать о частичной ``скрытой осведомлённости'' о ловушке; (8)~расовые эффекты существуют, но не образуют универсального паттерна: средние fallacy rate различаются менее чем на 1{,}5~п.п. между расовыми группами; (9)~две модели (\texttt{gemma-3-12b-it}, \texttt{grok-4.1-fast}) демонстрируют реверсивный эффект, при котором White-имена ассоциированы с несколько более высокой fallacy rate; (10)~профессиональные категории из сферы социального служения и образования (Social Services, Education, Healthcare) триггерят наибольший уровень предубеждения (35--40\% fallacy rate), тогда как финансовые и управленческие профессии --- наименьший (15--20\%); (11)~размеры эффекта Cohen's $d$ варьируют от пренебрежимого ($-0{,}112$ у \texttt{qwen3.5-397b}) до очень большого ($1{,}619$ у \texttt{claude-sonnet-4.6}); (12)~ячейки $n_{22}$ (двойная ошибка) и $n_{12}$ (обратная асимметрия) пусты для всех моделей, что подтверждает строго однонаправленную природу конъюнктивного предубеждения у LLM, отличающую их от человеческих испытуемых.

Сравнение с человеческой производительностью показало: (13)~наиболее подверженные предубеждению LLM (\texttt{qwen3-next-80b}, \texttt{gemma-3-27b-it}) демонстрируют fallacy rate (58--61\%), попадающий в диапазон между человеческими данными в probability- (85--89\%) и betting-форматах (~80\%), в то время как GPT-3.5 предыдущего поколения достигал человеческого потолка (90\%); (14)~три модели-лидера (\texttt{claude-sonnet-4.6}, \texttt{qwen3.5-397b-a17b}, \texttt{glm-4.7-flash}) с fallacy rate $\leq$ 0{,}8\% качественно превосходят людей даже в частотном формате (17--22\%), что указывает на принципиально иной уровень нормативного рассуждения; (15)~паттерн ошибки у LLM качественно отличается от человеческого: модели ошибаются исключительно в correlated-условии ($n_{22} = 0$), тогда как у людей нарушение конъюнкции наблюдается и при нейтральных стимулах, что свидетельствует о более «чистом» и стимульно-зависимом механизме эвристики репрезентативности у LLM.

